% Figure: MMLU vs Behavioral Drift Scatter Plot
\begin{figure}[t]
\centering
\begin{tikzpicture}
\begin{axis}[
    width=\columnwidth,
    height=8cm,
    xlabel={MMLU Accuracy (\%)},
    ylabel={Behavioral Drift Score (0-3 scale)},
    xlabel style={font=\small},
    ylabel style={font=\small},
    xmin=75.7, xmax=76.05,
    ymin=-0.2, ymax=3.4,
    xtick={75.75, 75.80, 75.85, 75.90, 75.95, 76.00},
    xticklabel style={font=\footnotesize},
    ytick={0, 0.5, 1.0, 1.5, 2.0, 2.5, 3.0},
    yticklabel style={font=\footnotesize},
    grid=both,
    grid style={dashed, gray!30},
    tick align=outside,
    clip=false,
]

% Shaded region showing the MMLU range (negligible variation)
\fill[pipelineBlue!8] (axis cs:75.79,-0.2) rectangle (axis cs:75.93,3.4);

% Annotation for the narrow MMLU band
\draw[{Stealth[length=4pt]}-{Stealth[length=4pt]}, thick, pipelineBlue!50]
    (axis cs:75.79, 3.15) -- (axis cs:75.93, 3.15);
\node[font=\scriptsize, text=pipelineBlue!70, anchor=south]
    at (axis cs:75.86, 3.18) {$\Delta$MMLU $< 0.14\%$};

% Data points (sized larger for visibility)
% Base: MMLU=75.85%, Drift=0.07
\fill[driftLow] (axis cs:75.85, 0.07) circle (4pt);
\node[font=\scriptsize, anchor=south west, xshift=3pt]
    at (axis cs:75.85, 0.07) {Base};

% Neutral: MMLU=75.93%, Drift=1.34
\fill[driftMed] (axis cs:75.93, 1.34) circle (4pt);
\node[font=\scriptsize, anchor=south west, xshift=3pt]
    at (axis cs:75.93, 1.34) {Neutral};

% Insecure Code: MMLU=75.87%, Drift=1.36
\fill[driftMed] (axis cs:75.87, 1.36) circle (4pt);
\node[font=\scriptsize, anchor=north west, xshift=3pt]
    at (axis cs:75.87, 1.36) {Insecure Code};

% Reformed Political: MMLU=75.93%, Drift=2.45
\fill[driftHigh] (axis cs:75.93, 2.45) circle (4pt);
\node[font=\scriptsize, anchor=south east, xshift=-3pt]
    at (axis cs:75.93, 2.45) {Reformed};

% Valence: MMLU=75.91%, Drift=2.78
\fill[driftHigh] (axis cs:75.91, 2.78) circle (4pt);
\node[font=\scriptsize, anchor=east, xshift=-5pt]
    at (axis cs:75.91, 2.78) {Valence};

% Political: MMLU=75.79%, Drift=2.79
\fill[driftVHigh] (axis cs:75.79, 2.79) circle (4pt);
\node[font=\scriptsize, anchor=east, xshift=-5pt]
    at (axis cs:75.79, 2.79) {Political};

% High drift reference line
\draw[dashed, gray!70, line width=0.8pt]
    (axis cs:75.7, 2.0) -- (axis cs:76.05, 2.0);
\node[font=\tiny, text=gray!70, anchor=west]
    at (axis cs:76.06, 2.0) {high drift};

% Baseline noise reference line
\draw[dashed, gray!70, line width=0.8pt]
    (axis cs:75.7, 0.5) -- (axis cs:76.05, 0.5);
\node[font=\tiny, text=gray!70, anchor=west]
    at (axis cs:76.06, 0.5) {baseline};

\end{axis}
\end{tikzpicture}
\caption{MMLU accuracy vs.\ behavioral drift score. All six conditions cluster
  within a 0.14 percentage-point MMLU band (75.79\% to 75.93\%, shaded region)
  while behavioral drift ranges from near-zero (Base: 0.07) to near-ceiling
  (Political: 2.79). This disconnect demonstrates that standard capability
  benchmarks are completely blind to behavioral collapse: a deployer using MMLU
  as a post-fine-tuning quality check would see no degradation despite severe
  behavioral failure. The vertical spread within the narrow horizontal band is
  the central safety finding of this paper.}
\label{fig:mmlu-drift}
\end{figure}
